\documentclass[leqno, a4paper, 12pt]{jsarticle}
\usepackage{amssymb}
\usepackage{amsthm}
\usepackage{amsmath}
\usepackage{comment}
\usepackage{url}

\usepackage{enumitem}
%\usepackage{showkeys}


%定理環境 thmはあまり使わずpropをなるべく使う。
%主張は基本的に証明の中で使い、そのため番号を付けない。
%注意環境を付けてみた。
\newtheorem{thm}{定理}
\newtheorem{prop}[thm]{命題}
\newtheorem{cor}[thm]{系}
\newtheorem{lem}[thm]{補題}
\newtheorem{df}[thm]{定義}
\newtheorem{exam}[thm]{例}
\newtheorem{fact}[thm]{事実}
\newtheorem*{claim}{主張}
\newtheorem*{cau}{注意}
\newtheorem*{rmk}{注意}
\renewcommand{\proofname}{証明}
%矢印たち。
%\toは\rightarrowと同じ。\getsは\leftarrowと同じ。同値は\iff
\newcommand{\To}{\Longrightarrow}
\newcommand{\Gets}{\LongLeftarrow}
%黒板太字
\newcommand{\nn}{\mathbb{N}}
\newcommand{\zz}{\mathbb{Z}}
\newcommand{\qq}{\mathbb{Q}}
\newcommand{\rr}{\mathbb{R}}
\newcommand{\cc}{\mathbb{C}}
%無限大
\newcommand{\mgn}{\infty}
%差集合は\setminusで統一する。等号の否定は\neq
%真部分集合は\subsetneqq　偏微分は\partial
%ディスプレイスタイル。\dfracでディスプレイ分数
\newcommand{\dpst}{\displaystyle}

\title{論文メモ
\large{\\--- カントール集合？ ---}
}
\author{ナンブ　キトラ}
\date{\today}
\begin{document}
\maketitle

本棚を整理したいので，
溜まっている論文のメモを書いて未来への手紙とすることにする．


\section{本文}
まず論文
\cite{garity2007cantor}
はカントール集合に関する未解決問題の
解説である．
未解決問題を参照したいときはこれを読むといいだろう．


\subsection{カントール集合の形}

論文
\cite{MR1203278}
は何やら自己相似集合がいい感じだとカントール集合といい感じでヘルダー同値ということを言っている．


論文
\cite{MR0961514}
はアフィンなカントール集合を分類している気がするがよくわからず．この論文J. Diff. Geom. で出版されてるのが面白いと思う．

論文
\cite{bellissard2012bi}
はカントール集合をユークリッド空間に埋め込む話をしている．


\subsection{カントール集合と測度}
測度に関する話は以下の論文である．

論文
\cite{MR3153999}
では自己相似ではないカントール集合の測度を計算しているようだ．


論文
\cite{MR4487615}
では三進カントール集合
を$\mathcal{K}$
として$P\colon \rr^{2}\to \rr$
を$P(x, y)=xy$
としたとき
\[
\left|\mathcal{L}^{1}(P(\mathcal{K}\times \mathcal{K}))-
\frac{91782451}{113374080}\right|<\frac{1}{10^{6}}
\]
であることを証明している．
なんかやたら具体的な数字が出てきたので
面白がって印刷したんだと思う．

\subsection{野生のカントール集合}
Wild Cantor set を野生の...と訳すのは
ちょっと間違いなのだが面白のでそう書いておく．
多分鎖に繋がれていないとか人に慣れてないとかそういうニュアンスだと思う．tameの反対．
なんかよくわかんないので
論文を記録に残すだけにする．
先述の\cite{garity2007cantor}
も野生のカントール集合の話である．

\cite{MR0375323}

\cite{MR3303046}

\cite{MR0896874}

\cite{MR0234438}


\cite{MR0580587}


\cite{MR0221487}

\subsection{野生のcell}
なんか混ざってた．

\cite{MR0027512}

\cite{MR0178460}


\subsection{その他}
論文
\cite{MR3397281}
は力学系の話だと思うがわからず．

論文
\cite{MR2291714}
は
カントール集合の
自己同相群を調べている．

論文
\cite{MR3904605}
では$[0, 1]$の元が全て$u=x^{2}y$, ここで
$x, y$はカントール集合の元
で書かれるということを証明している．
他にも$x/y$
や$xy$
のようなものについて調べている．
論文
\cite{MR4487615}
と関係があるだろう．

%READ!
%myplain is the same to amsplain style. 
\nocite{*}
\bibliographystyle{myplaindoidoi}
\bibliography{cantorset.bib}



\end{document}